% Online Appendix of the AEJ:EP version (support_global_redistr_AEJEP.tex). Anonymized.
% Section A gathers material moved from the main text of support_global_redistr.tex; the other sections are app.tex, unchanged.

\documentclass[12pt,english]{article}
\usepackage[utf8]{inputenc}
\usepackage[T1]{fontenc}
\usepackage{tgpagella}
\usepackage{mathpazo}
\usepackage[margin=1.5in]{geometry}
\usepackage{fullpage} % page layout as in support_global_redistr.tex, for which appendix figures are sized
\usepackage{hwemoji}
\usepackage[round,sort&compress]{natbib}
\usepackage{url}
\usepackage{xr-hyper}
\usepackage{xcolor}
\usepackage{hyperref}
\hypersetup{colorlinks=true, urlcolor=purple, linkcolor=blue, citecolor=violet}
\externaldocument{support_global_redistr_AEJEP}
\usepackage{bm}
\usepackage{indentfirst}
\setcitestyle{aysep={}}
\usepackage[leqno]{amsmath}
\usepackage{amssymb}
\usepackage{eurosym}
\usepackage{babel}
\usepackage{graphicx}
\usepackage{caption}
\usepackage{subcaption}
\usepackage{float}
\usepackage{enumitem}
\usepackage{setspace}
\usepackage{gensymb}
\usepackage{textcomp}
\usepackage{tabularx}
\usepackage{booktabs}
\usepackage{makecell}
\usepackage{paracol}
\usepackage{morefloats}
\usepackage[export]{adjustbox}
\usepackage{pdflscape}
\usepackage{longtable}
\renewcommand{\floatpagefraction}{.99}
\extrafloats{300}
\usepackage{lineno} % lineno also prevents an output-routine dead cycle with the too-large conjoint-by-vote figure of app.tex
\modulolinenumbers[5]
\linenumbers

\renewcommand{\thetable}{S\arabic{table}}
\renewcommand{\thefigure}{S\arabic{figure}}

\begin{document}

\begin{center}
  {\LARGE \bfseries Online Appendix}\\[6pt]
  {\Large Public Acceptance of International Redistribution in High-Income Countries}
\end{center}

\tableofcontents
\clearpage

\appendix

%%%%%%%%%%%%%%%%%%%%%%%%%%%%%%%%%%%%%%%%%%%%%%%%%%%%%%%%%%%%%%%%%%%%%%%%%%%%%%%%
\section{Additional Material on the Main Analyses}\label{app:moved}
%%%%%%%%%%%%%%%%%%%%%%%%%%%%%%%%%%%%%%%%%%%%%%%%%%%%%%%%%%%%%%%%%%%%%%%%%%%%%%%%

\subsection{Country Coverage}

\begin{figure}[h!]
    \caption[Country coverage]{Country coverage of the survey.}\label{fig:country_coverage}
    \makebox[\textwidth][c]{\includegraphics[width=\textwidth]{../figures/maps_participation/country_coverage_curtailed.png}}
\end{figure}

\subsection{Open-Ended Questions}\label{app:fields}

At the beginning of the survey, respondents are randomly assigned one of four open-ended questions: their main concerns, their needs or wishes, an issue important to them but neglected in public debate, or the greatest injustice of all. The questions are deliberately broad and vague to let respondents express their top-of-mind considerations without priming.

\paragraph{Classification.} I automatically translated each answer into English.\footnote{I used \textit{onlinedoctranslator.com}, which is powered by \textit{Google Translate}.} I then used AI and my own reading of a few hundred answers to identify the most common concepts, from which I selected 27 categories. Next, I classified each answer into one or more of these categories, both manually (Figures \ref{fig:field_manual}--\ref{fig:injustice_field}) and automatically using AI (Figure \ref{fig:field_gpt}). Finally, I manually defined a list of 47 (disjunctions of) keywords and used it to classify all responses automatically (Appendix~\ref{subsec:keywords}; Figure~\ref{fig:field_keyword}).

The three classification methods yield consistent results but differ in accuracy. The keyword classification allows for an exact and reproducible search, while the AI classification is not limited to specific words and captures more matching responses. The manual classification seems to provide the most accurate results, with a number of matches generally between those of the other two methods. For example, for the \textit{injustice} question, 1.2\% of answers match the keywords for \textit{global inequality}, ChatGPT identifies this category in 7.5\% of answers, and my manual coding in 3.2\%.\footnote{The keyword matching searches the regular expression \texttt{global poverty\allowbreak|global inequal\allowbreak|hunger|drinking water|starv}, ignoring case. The automatic and manual classifications are based on the category definition ``Inequality at the international level / Hunger or poverty in poor countries''.} The AI incorrectly classifies unspecific answers like ``poverty'' in this category,\footnote{Out of the 47 one-word answers ``poverty'', (the zero-shot prompt passed to) ChatGPT-4.1 coded 42 as \textit{global inequality}, illustrating the lack of consistency of this classifier.} while the keyword search misses answers like ``inequality among humans''. I therefore use the manual classification as the benchmark and the two other methods as robustness checks.

\paragraph{Results.} Word clouds (Figure \ref{fig:wordcloud}) visualize the most common concepts. The most frequent \textit{concerns} or \textit{wishes} relate to purchasing power, with concepts such as ``money'', ``inflation'', the ``cost of living'', or ``financial stability'' appearing in 31\% of these answers. Within countries, the share of people concerned with money decreases with income, from 35\% in the bottom decile to 22\% in the top decile.\footnote{At the country level, the concern for money is significantly correlated with inequality: an additional point in the Gini index is associated with 0.8~p.p.\ more respondents concerned with money. The concern for money is higher in richer countries, but the correlation vanishes once one controls for the Gini.} The next most frequent \textit{concerns} are health (13\%), far-right governments or related concepts such as ``Trump'' or ``trade tariffs'' (10\%), and war (9\%). Most \textit{wishes} are personal: after money come one's own or one's relatives' health (21\%) and peace of mind (10\%). The predominance of material considerations is consistent with previous studies \citep{wlezien_salience_2005,singer_who_2011}.

Asked about the greatest \textit{injustice}, 19\% of answers relate to ``inequality'' or ``poverty'' (28\% in Europe but only 9\% in the US). Only 11\% of these answers specify a geographical scope: 2\% mention their own country versus 10\% the global level (or Global South issues such as ``clean water'' or ``starvation''). Italians, Poles, and Spaniards are the most likely to mention global inequality or poverty, while Japanese and Russian respondents are the least likely. The next most common answers relate to discrimination (9\%), violence or wrongful convictions (9\%), or the welfare state (8\%), with respondents criticizing either the lack of public services or excessive welfare given to the undeserving.

Asked about ``an issue important to [them] but neglected in the public debate'', 21\% of respondents cannot name one. The most frequently mentioned concepts are public services (12\%), the cost of living (10\%), health (9\%), ageing (6\%), and the environment (6\%); immigration is mentioned in 5\% of answers. That the most frequently mentioned topics are already well publicized suggests that public debate reflects or shapes what people have in mind. Appendix~\ref{subsec:slant} lists topics that are uniquely prevalent in each country.

The topics mentioned vary with sociodemographic characteristics. For example, a respondent who mentions \textit{immigration} is 3.5 times more likely to vote for the far right (correlation of .17), and one who mentions \textit{old age} is twice as likely to be 65 or older (.13). Other strong correlations are between \textit{criticizing the far right} and voting for the left (.16), \textit{health} and age (.11), \textit{employment} and being unemployed (.09), \textit{animals} and extending one's moral circle to sentient beings (.09), \textit{education} and being a student (.09), \textit{the environment} and voting for the left (.08), and \textit{money} and income ($-$.08).

\begin{figure}[h!]
  \caption[Word cloud of open-ended field, per variant]{Most common concepts in open-ended fields. (Questions~\ref{q:concerns_field}--\ref{q:issue_field})} \label{fig:wordcloud}
  \begin{subfigure}{.48\textwidth}
    \caption[]{``What are your main concerns these days?''}
    \includegraphics[height=.35\textheight]{../figures/all/concerns_field_en.pdf}
  \end{subfigure} \quad
  \begin{subfigure}{.48\textwidth}
    \caption[]{``What are your needs or wishes?''}
    \includegraphics[height=.35\textheight]{../figures/all/wish_field_en.pdf}
  \end{subfigure}
  \begin{subfigure}{.48\textwidth}
    \caption[]{``What according to you is the greatest injustice of all?''}
    \includegraphics[height=.35\textheight]{../figures/all/injustice_field_en.pdf}
  \end{subfigure} \quad
  \begin{subfigure}{.48\textwidth}
    \caption[]{``Can you name an issue that is important to you but is neglected in the public debate?''}
    \includegraphics[height=.35\textheight]{../figures/all/issue_field_en.pdf}
  \end{subfigure}
\end{figure}

\subsection{Maps of the International Climate Scheme}

\begin{figure}[h!]
\caption[Example maps of the International Climate Scheme]{Example maps of the International Climate Scheme question. The respondent's own country is striped to denote its potential participation. (Question~\ref{q:ics_support}).} \label{fig:ics_maps}
\begin{subfigure}{.49\textwidth}
  \caption[]{\textit{Low} variant for the US.}
  \includegraphics[height=.49\textwidth]{../figures/maps_participation/GCS_low_EN.png}
\end{subfigure}
\begin{subfigure}{.49\textwidth}
  \caption[]{\textit{High color} variant for EU countries.}
  \includegraphics[height=.49\textwidth]{../figures/maps_participation/GCS_high_color_EU.pdf}
\end{subfigure}
\end{figure}

\subsection{Preferred Channels for Transferring Resources to LICs}\label{subsec:transfer_how}

Asked to evaluate ways of transferring resources to reduce poverty in LICs on a 4-point scale, the most preferred option in every country is ``Cash transfers to parents (child allowances), to the disabled and to the elderly'', with 16\% selecting it as the \textit{Best way} overall, and 49\% as a \textit{Right way} or \textit{Best way} (Figures \ref{fig:transfer_how}, \ref{fig:transfer_how_above_one}--\ref{fig:transfer_how_negative}).\footnote{It would be more relevant to ask this question to the population of LICs. \cite{carson_publics_2025} survey seven Global South countries and find that people favor using carbon tax revenues for public services rather than cash transfers.} ``Unconditional transfers to the national government'' is the only option seen as a \textit{Wrong way} by a majority, but this share falls from 51\% to 21\% (making it the third most supported option out of seven) when ``transfers to the national government [are] conditioned on the use of funds for poverty reduction programs''. ``Unconditional cash transfers to each household'' are controversial: they are the second most chosen \textit{Best way}, yet 33\% view them as a \textit{Wrong way}. ``Transfers to public development aid agencies which then finance suitable projects'' are uncontroversial, with only 16\% considering them a \textit{Wrong way} and 37\% a \textit{Right} or \textit{Best way}.

\begin{figure}[h!]
    \caption[\textit{Right} or \textit{Best way} to transfer resources to LICs]{``How do you evaluate each of these channels to transfer resources to reduce poverty in LICs?''\\ Percentage of \textit{Right} or \textit{Best way} (other options: \textit{Wrong} or \textit{Acceptable way}). (Question~\ref{q:transfer_how}).}\label{fig:transfer_how}
    \makebox[\textwidth][c]{\includegraphics[width=\textwidth]{../figures/country_comparison/transfer_how_positive.pdf}}
\end{figure}

\subsection{Moral Circle}

\begin{figure}[h!]
    \caption[Moral circle]{``Which group of people do you advocate for when you vote?'' In Russia and Saudi Arabia, the question read: ``Which group do you advocate for? For example, if you were the richest person on Earth, which group would you predominantly help with your money?'' (Question~\ref{q:group_defended}).}\label{fig:group_defended}
    \makebox[\textwidth][c]{\includegraphics[width=.9\textwidth]{../figures/all/group_defended_nolabel.pdf}}
\end{figure}

Universalists are fewer in Japan (30\%) and constitute a majority in Europe (50\%) and Saudi Arabia (57\%) (Figures \ref{fig:group_defended} and \ref{fig:group_defended_all}). Among those who lean to the left, 59\% are universalists, compared to 32\% on the center-right or far-right.

Following \cite{enke_universalism_2024}, I construct an alternative measure of universalism based on the vocabulary used in open-ended fields.\footnote{I use the Moral Foundation Dictionary (MFD) 2.0 \citep{frimer_moral_2019} and define \textit{universalism} as the number of occurrences of \textit{care} or \textit{fairness} words minus the number of \textit{loyalty} or \textit{authority} words. An alternative definition based on the extended MFD \citep{hopp_extended_2021}, which uses weights rather than dummy variables, is even less (though still significantly) correlated with latent support for global redistribution, at .03.} Although latent support for global redistribution is significantly correlated with this measure, the correlation is only .05, lower than the correlation with the manual, keyword, or AI classifications of an answer as relating to inequality or poverty (.09, .08, and .06), and much lower than the correlation with the moral-circle measure of universalism (.37). The correlation between the two measures of universalism is only .03, which shows that indicators labeled as ``universalism'' by different authors may not capture the same dimension.

\subsection{Custom Global Income Redistribution}\label{subsec:custom_redistr}

The last task of the questionnaire lets respondents manipulate the shape of the global income distribution.\footnote{Appendix \ref{subsec:country_features} details how I obtained the world distribution of PPP incomes.} The question text included the following instructions:
\begin{quote}``Below you will find a graph of the world distribution of after-tax income and three sliders that vary it. The current distribution is in red, and your custom one is in green. The first two sliders control the proportion of winners and the proportion of losers, among all humans. The third slider controls the degree of redistribution from the richest to the poorest. If you do not want new policies to reduce global inequality, you can set the third slider to zero.''\end{quote}
Appendix \ref{app:algo} describes the algorithm translating slider positions into a redistribution. Figure \ref{fig:custom_redistr_question} displays what respondents see below the instructions: the interactive graph and a table summarizing how their custom redistribution would affect five example income levels, including their own. To mitigate anchoring, the sliders are initialized in one of two random positions: either 60\% of winners, 20\% of losers, and a degree of redistribution of 2 out of 10, or 40\%, 10\%, and 7/10.\footnote{To test for anchoring, I regress responses on the sliders' initial positions and define the anchoring effect as the effect size relative to the difference between the two initial positions. It is always significant: 36\% for the share of winners, 57\% for the share of losers, and 42\% for the degree of redistribution. Responses converge to a point between the two anchors.} Given the complexity of the task and its inconvenience on mobile devices, respondents are explicitly allowed to skip it.

\begin{figure}[h!]
   \caption[Custom redistribution task screenshot]{Custom global redistribution: screenshot of the bottom of the page. (Question~\ref{q:custom_redistr}).}\label{fig:custom_redistr_question}
    \makebox[\textwidth][c]{\includegraphics[height=.8\textheight]{../figures/questionnaire/survey_custom_redistr_bottom.png}}
\end{figure}

Overall, 56\% are satisfied with their custom redistribution and 43\% skip the task. Non-response is spread relatively evenly across the population: the share of satisfied respondents is 52\% among non-voters, 54\% among center-right or right-wing voters, 57\% among far-right voters, and 61\% among left-wing voters, and it ranges from 49\% among people without a high school diploma to 57\% among those with a post-secondary diploma.

Figure \ref{fig:custom_redistr_median} shows the median preferred redistribution among satisfied respondents, obtained by setting the sliders at their median values: 49\% of winners, 18\% of losers, and a degree of redistribution of 5/10, resulting in a transfer of 5.4\% of world income and a minimum income of \$287 per month (PPP). 48\% choose to lose from their custom redistribution, while only 9\% choose to win; the median satisfied respondent neither wins nor loses. 10\% of satisfied respondents opt for the status quo. The average preferred redistribution, obtained by pointwise averaging of custom curves (Figure \ref{fig:custom_redistr_mean}), transfers 5.4\% of world income from the top 27\% to the bottom 73\% and entails a minimum income of \$247 per month. At the top of the distribution, it can be achieved with a 7\% marginal income tax rate above \$25,000 and a 16\% rate above \$40,000 per year (PPP; Figure \ref{fig:tax_radical_redistr}). The average (resp.\ median) preferred redistribution would reduce the world Gini index from .67 to .59 (resp.\ .58) and the D9/D1 ratio from 42 to 11 (resp.\ 10), and would increase the bottom 50\% share from 8\% to 13\%. Heterogeneity across countries is limited (Figures~\ref{fig:custom_redistr_satisfied_mean}--\ref{fig:custom_redistr_satisfied_median}), while individual choices are more dispersed, although most respondents favor transfers of 4\% to 5\% of world income and a minimum income of \$150 to \$350 per month (Figures~\ref{fig:custom_redistr_winners_agg}--\ref{fig:custom_redistr_transfer_ceiling}).

\cite{fabre_french_2022} applied the same method to French preferences for national redistribution and tested support for the median and average preferred redistributions on a separate sample: excluding the 22\% to 24\% of non-respondents, 51\% accepted the average redistribution and 67\% the median one. Whether these results replicate for global redistribution remains to be tested.

\clearpage
\input{app}

\clearpage
\singlespacing
\bibliographystyle{aea}
\bibliography{global_tax_attitudes,aejep_extra}

\end{document}
